\section{Fountain 码与无速率传输}
\label{sec:04-Discrete-Mathematics/05-Coding-and-Polynomial-Methods/05}

一份几十 GB 的模型权重要推给上万台机器。有的机器在同一机房，丢包率千分之几；有的挂在跨洋链路末端，一阵抖动就少掉一大块。发送方事先并不知道谁会丢多少。

固定码率的分组码在这里很别扭：冗余要在发包之前就定死。按 \(10\%\) 的丢包率配冗余，实际只丢 \(1\%\) 的那批机器白白多收了九成的校验数据；而那台丢了 \(40\%\) 的机器，一个符号也恢复不出来，整份数据作废。自动重传对一对一通信没问题，一对上万就成了灾难——每台机器缺的包各不相同，发送方要为每一台维护一份反馈状态，然后把几乎所有包都重发一遍。

Fountain 码换了个提法。

\begin{quote}
编码器不预先决定要发多少，它持续产生新的编码包；任何接收者只要收集到略多于原始包数的不同编码包，就能独立还原全部数据，无需告诉发送方自己收到了哪些。
\end{quote}

``无速率''指的是码长不必事先固定，与传输速度无关。

\begin{center}
\begin{tikzpicture}[x=1cm,y=1cm,font=\small,>=stealth]
  \draw[black!55,fill=blue!10] (0,2.5) rectangle (1.5,3.9);
  \node[font=\footnotesize] at (0.75,3.2) {编码器};
  \node[below,font=\footnotesize,black!70] at (0.75,2.44) {\(x_1,\ldots,x_k\)};
  \draw[->,black!45] (1.6,3.4) -- (2.4,3.7);
  \draw[->,black!45] (1.6,3.2) -- (2.5,2.9);
  \draw[->,black!45] (1.6,3.0) -- (2.3,2.25);
  \foreach \x/\y in {2.9/3.5,3.6/3.0,4.3/3.6,5.0/2.9,5.7/3.5,6.4/2.9,7.1/3.5,7.6/2.7,3.1/2.2,4.0/1.9,4.9/2.4,5.6/1.9,6.5/2.2,7.2/1.8}{
    \fill[blue!55] (\x,\y) circle (1.9pt);
  }
  \node[right,font=\footnotesize,black!70] at (7.9,2.8) {\(\cdots\) 不设终点};
  \draw[black!65,thick] (2.6,1.4) -- (2.6,0.6) -- (4.0,0.6) -- (4.0,1.4);
  \draw[black!65,thick] (5.6,1.4) -- (5.6,0.6) -- (7.0,0.6) -- (7.0,1.4);
  \foreach \x/\y in {2.9/0.85,3.3/0.8,3.7/0.9,3.1/1.15}{\fill[blue!70] (\x,\y) circle (1.9pt);}
  \foreach \x/\y in {5.9/0.85,6.3/0.8,6.7/0.9,6.1/1.15}{\fill[red!60] (\x,\y) circle (1.9pt);}
  \node[below,font=\footnotesize] at (3.3,0.52) {接收者甲};
  \node[below,font=\footnotesize] at (6.3,0.52) {接收者乙};
  \draw[->,dashed,black!35] (3.1,2.2) -- (3.2,1.5);
  \draw[->,dashed,black!35] (4.0,1.9) -- (3.8,1.5);
  \draw[->,dashed,black!35] (5.6,1.9) -- (5.9,1.5);
  \draw[->,dashed,black!35] (6.5,2.2) -- (6.5,1.5);
\end{tikzpicture}

\emph{两个接收者接到的包互不相同，谁也没有向上游报告过自己缺什么。各自接够略多于 \(k\) 个就能独立还原——要付出的就藏在``略多于''三个字里，以及一个概率性的、而非确定性的承诺。}
\end{center}

\subsection{每个编码包是一条随机方程}

把原始数据切成 \(k\) 个等长的包 \(x_1,\ldots,x_k\)。生成一个编码包时，先按某个度分布抽一个度数 \(d\)，再随机挑 \(d\) 个原始包异或起来：
\[
y=x_{i_1}\oplus\cdots\oplus x_{i_d}.
\]
包头写明参与异或的下标，或者只写一个能重现这些下标的随机种子。于是接收者手里的每个包都是关于未知量 \(x_i\) 的一条 \(\mathbb F_2\) 上的线性方程，而它们的地位完全平等——没有哪个包是``第几号校验包''，也没有哪个包不可或缺。

方程攒够且线性无关，做 Gauss 消元就能解出全部未知量。但稠密矩阵消元的开销随 \(k\) 的三次方增长，几万个包时不可接受。LT 码的贡献是把度分布设计好，让绝大部分求解可以由一个近乎线性时间的过程完成。

\subsection{剥离要靠度为一的包起步}

那个过程与上一节擦除信道上的连锁剥离是同一套动作。收到一个度为一的包 \(y=x_3\)，\(x_3\) 立刻确定；把它从所有含它的方程里异或掉，那些方程的度数各减一；某条方程若因此降到度一，又解出一个新的原始包。骨牌一路推下去，直到全部解开或者中途卡住。

度分布因此夹在两个要求之间。度都很大，开局就没有度一的包，骨牌根本推不倒第一张；度都很小，包与包反复覆盖同一小撮原始包，剩下的那些始终没人碰。LT 码的鲁棒孤子分布就是为了在整个剥离过程中维持一个数量适中的度一候选池——不多，多了浪费；不少，少了链条断掉。Raptor 码更进一步，先用一层高码率预编码把原始包扩充一遍，容忍剥离漏掉少量包，再在扩充结果上跑简单的 LT 编码，把开销和复杂度一起压下来。

剥离卡住不等于信息不足。有时矩阵仍然满秩，只是这套贪心规则用尽了；这时对剩下的小系统做一次 Gauss 消元，或者用失活译码把少数变量暂时挂起、继续剥离其余部分，往往就能走完。判断``算法停了''和判断``方程不够''是两件事，实现里混淆这两者会导致无谓的重传。

\subsection{那个``略多于''是随机性的定价}

随机抽出来的方程有可能重复，也可能线性相关，所以恰好收到 \(k\) 个包并不保证满秩。实际方案要收到大约 \(k(1+\varepsilon)\) 个，\(\varepsilon\) 的大小取决于度分布、块长和可接受的失败概率。

这一点值得与上一节的 Reed--Solomon 摆在一起看。Reed--Solomon 给的是确定性保证：任意 \(k\) 个正确的求值点，一定唯一确定那个多项式，没有例外，也不存在``运气不好''。Fountain 码给的是概率保证：从设计好的分布里随机取略多于 \(k\) 个包，以高概率能还原。把它说成``任意 \(k\) 个编码包都能恢复''是不对的——那正是它为了摆脱固定码率而放弃的东西。

换来的是别的：编码器不必知道信道，不必收反馈，不必为不同接收者准备不同的冗余，想发多久发多久。在一对多、无反馈、擦除率未知的场景里，这笔交换很划算；在一对一、丢包率稳定、又要求确定性容错的场景里（比如那个 6 加 3 的存储条带），它不如 Reed--Solomon。选码的第一步从来不是比较性能表，而是先说清信道是什么样、要的是哪一类保证。

\subsection{能纠丢包，不等于能识破篡改}

Fountain 码假定接收者清楚哪些包没到，并且到了的包内容都是对的。若有一个包被悄悄改了几个比特而接收者不知情，这个错误会在剥离过程中沿着异或链扩散到许多原始包上——一个坏包毁掉一大片，比丢了它还糟。

所以真实系统总要另配校验和、消息认证码或签名，先确认包的完整性，再交给译码器。纠擦除、检测随机错误、抵御恶意篡改是三件不同的事，靠三种不同的机制；一种编码能从丢包中恢复数据，不代表它顺带提供了真实性。这也是上一节末尾那句话的另一面：与其让纠错机制去猜哪一份数据可疑，不如让每一份数据自证清白。

\subsection{整章留下的判断}

从 Reed--Solomon 到 LDPC 再到 Fountain，工具从多项式换到稀疏图再换到随机异或，导读那幅图始终没变：把 \(q^k\) 个消息稀疏地撒进一个 \(q^n\) 个点的空间，让它们彼此隔开，噪声推不出各自的势力范围就还能认回去。区别只在于用什么关系把码字绑住，以及绑法允许什么样的译码算法。

选择时真正要问的是两件事。信道给的是错误、擦除还是软噪声——位置已知与否，价钱差一倍；要的是确定性保证还是高概率保证——前者写进合同，后者写进服务等级目标。这两件事定下来，码族的取舍范围基本就收敛了；反过来，先挑码再论证信道，多半会在某个尾部场景上失效。

还有一条边界，本章自始至终没有触碰：给定信道，可靠传输的码率上限究竟是多少？这里的每条结论都是关于某个具体构造能做到什么，而不是任何构造都越不过什么。那个上限由信道容量给出，是\hyperref[ch:058]{``有噪信道编码：噪声中为什么仍能可靠通信''}的主题——LDPC 与 Turbo 码所逼近的``极限''，正是在那里被定义出来的。
